\chapter{Evaluation}
For a project like AirScout it is essential to determine a fitting outlier-detection method. Thus it's important to compare the two available methods for outlier detection mentioned in this chapter ~\ref{title:outlier}.

\boldtitle{Possible methods}
The two possible strategies are:
\begin{itemize}
    \item Z-Score method
    \item IQR-rule
\end{itemize}

\boldtitle{Evaluation criteria and weighting}
Before it is possible to rank the strategies it is first necessary to define evaluation criteria:
\begin{table}[H]
    \centering
    \begin{tabular}{|p{4.5cm}|p{4.5cm}|p{5cm}|}
    \hline
   Criteria & Description & Scale \\
   \hline
   K1: Robustness against outliers & How strongly is the method affected by extreme values in the dataset? & 1 = very sensitive, 5 = very robust\\
   \hline
   K2: Assumptions about data distribution & How dependent is the method on the underlying data distribution (e.g. normal dsitribution)? & 1 = very strict assumptions, 5 = no assumptions\\
   \hline
   K3: Computation cost (local) & How much computational effort is required to calculate the method? & 1 = very high, 5 = very low\\
   \hline
   K4: Stability for small datasets & How reliable is the method when only a small amount of data is available? & 1 = very unstable, 5 = very unstable\\
   \hline
   K5: Implementation complexity & How complex is the implementation of the method? & 1 = complex, 5 = very simple\\
   \hline
    \end{tabular}
    \caption{Defined criteria}
    \label{tab:criteria}
\end{table}

Finally a weight needs to be added to each of the criteria:
\begin{table}[H]
    \centering
    \begin{tabular}{|p{8cm}|p{2cm}|}
    \hline
    Criteria & Weight\\
    \hline
       K1: Robustness against outliers & 30\%\\
       \hline 
       K2: Assumptions about data distribution & 25\%\\
       \hline
       K3: Computation cost (local) & 15\%\\
       \hline
       K4: Stability for small datasets & 15\%\\
       \hline
       K5: Implementation complexity & 15\%\\
       \hline
    \end{tabular}
    \caption{Defined weights}
    \label{tab:weights}
\end{table}

K1 receives a weight of 30\% because robustness against outliers is essential for reliable results in noisy, real-world sensor data. K2 is also weighted at 25\% since incorrect assumptions about data distribution can significantly affect the validity of the method. K3 receives 15\% because computational effort is relevant for embedded systems but not the primary constraint. K4 is assigned 15\% as stability with small datasets reliability, but is less critical than robustness and correctness. K5 gets 15\% since implementation effort matters, but has less impact than the statistical properties of the method.

\boldtitle{Ranking the options}
On the now defined criteria it is possible to give them a score and ultimately rank them:
\begin{table}[H]
    \centering
    \begin{tabular}{|c|c|c|c|c|}
    \hline
    Criteria & Weight & Z-Score & IQR-Rule\\
    \hline
    K1 & 30\% & 3 & 5\\
    K2 & 25\% & 4 & 5\\
    K3 & 15\% & 5 & 2\\
    K4 & 15\% & 5 & 2\\
    K5 & 15\% & 5 & 4\\
    \hline
    Total & & 4.15 & 3.95\\
    \hline
    \end{tabular}
    \caption{Ranking the options}
    \label{tab:placeholder}
\end{table}

The strategy with the best score is the Z-Score method. Thus AirScout devices will filter outliers based on the Z-Score method.